\section{Ablation and retained component comparisons}\label{sec:ablation}
A causal ablation would intervene on one implementation feature
while preserving the same problem, inputs, acceptance criteria and
other numerical choices. The retained evidence does not perform
every such intervention. Table~\ref{tab:ablation} instead states
which component comparisons actually exist and what they can
support. This distinction prevents a required analysis section
from being filled with an invented experiment or an unsupported
causal explanation for the primary timing outcome.

\begin{table}[htbp]\centering\small
\caption{Scope of retained component comparisons and unperformed causal interventions. Each row distinguishes recorded numerical evidence from a stronger attribution that the existing experimental design does not establish.}\label{tab:ablation}
\begin{tabularx}{\textwidth}{@{}lXX@{}}\toprule
Component & Existing comparison & Inference limit\\\midrule
Representation & Native, compact epigraph and row-split policies & Complete policies differ; no isolated causal factorization effect\\
Workspace reuse & Common C12 cold/warm observations & No independent intervention on every initial iterate, cache and active set\\
Conditioning & Fixed-within-seed diagonal condition sweep & Finite robustness grid, not arbitrary physical scaling\\
Directional policy & Same-input equivalent teaching formulations & Correctness illustration, not timed performance or field benefit\\
Threshold implementation & Native exact prefix operation & No alternative sorting implementation timed in isolation\\
Acceptance overhead & Included original-space checks & No post-hoc removal to improve the registered endpoint\\\bottomrule
\end{tabularx}
\end{table}

The formulation comparison retains exact objectives. Row splitting
changes row count and scaling while using a published grouped
projection backend; a compact epigraph changes variable structure
and can interact differently with sparse ordering and numerical
preparation. The returned time distributions therefore compare
complete reference policies. They do not isolate one matrix-size
effect or prove that a threshold sort explains every difference.
The unequal compatibility and failure counts in
Figure~\ref{fig:coupled-formulations} further prevent reading its
unmatched boxes as a causal ranking on one common successful subset.

Cold/warm reuse provides a different component view. All methods
receive the same oracle-driven C12 states, so differences are not
caused by independently divergent plant trajectories. However,
warm mode may retain factors, active sets and iterates together,
while setup accounting also changes. The stream has one observation
per frame and remains descriptive. It cannot assign the entire
reduction to any one retained object or supply the primary
confidence interval. Online active-set literature already uses
parametric structure~\cite{qpOASES}; the observation is not evidence
that reuse was introduced by this implementation.

The conditioning sweep changes diagonal curvature within its
registered seed controls. It supports a bounded statement about
objective/hard-error behavior on that grid. The dimensional
teaching checks test a different invariant: reciprocal row and
directional-scale conversion should preserve physical meaning.
Changing a directional scale alone changes the objective's
allocation policy. Neither comparison is a substitute for an
unperformed randomized intervention on threshold arithmetic or
for a new registered timing study. The fixed primary outcomes
therefore remain the inferential conclusions, while these
component analyses explain what future mechanism-specific
experiments would need to hold constant.
